\documentclass[11pt]{article}

\usepackage[a4paper,margin=1in]{geometry}
\usepackage{amsmath}
\usepackage{amssymb}
\usepackage{booktabs}
\usepackage{tabularx}
\usepackage[hidelinks]{hyperref}
\usepackage[authoryear,round]{natbib}
\usepackage{siunitx}
\sisetup{group-separator={,},group-minimum-digits=4}

\newcommand{\Hnorm}{1-H_{\mathrm{norm}}}

\title{Controlled Evaluation of Class-Imbalance Mitigation\\
Across Histopathology Patch Classification and WSI-Level MIL}
\author{Yannik Queisler}
\date{Experimental protocol, \today}

\begin{document}
\maketitle

\begin{abstract}
\noindent This report prespecifies a controlled benchmark of class-imbalance mitigation in computational pathology. Four datasets---CAMELYON16, TCGA-UT, BRACS, and PANDA---provide patch-classification and whole-slide multiple-instance-learning (MIL) regimes with frozen Virchow2 features. Patient-disjoint splits are created before imbalance is manipulated, and only training support changes: a full natural anchor is accompanied by size-matched balanced, moderate ($\rho=10$), and severe ($\rho=100$ where feasible) conditions. Validation and test patients, held-out prevalence, training-set size, model family, and evidence budgets remain fixed. The primary endpoint is balanced accuracy. Mitigation is evaluated as recovery of the deficit between size-matched balanced and imbalanced cross-entropy, conditional on evidence that such a deficit exists. Secondary analyses address classwise discrimination, calibration, computational cost, independent feature-space separability, effective sample size, and the interaction between slide imbalance and informative-instance sparsity in CAMELYON16. The document contains the complete design and statistical analysis plan but no experimental findings.
\end{abstract}

\setcounter{tocdepth}{2}
\tableofcontents
\clearpage

\section{Introduction}
Class imbalance changes the empirical objective faced by a classifier: frequent classes supply more optimization terms, while rare classes may contribute too little evidence to establish stable decision boundaries. In pathology, the problem occurs at two distinct prediction units. Patch classification uses directly labelled crops, whereas WSI-level MIL predicts from bags whose instances inherit only a weak slide label. The same slide cohort can therefore be nearly balanced between diagnoses while containing very sparse class-informative tissue. Aggregate accuracy does not isolate either failure mode, motivating class-balanced discrimination and probability-quality evaluation \citep{saini2023tacklingClass,zhang2023deepLong,mosquera2024classImbalance}.

Native-dataset comparisons do not identify an imbalance effect by themselves. They simultaneously vary task, class separability, training support, label granularity, and evaluation prevalence. This protocol instead changes class allocation only in the training partition and retains identical natural validation and test cohorts. A size-matched balanced CE condition estimates attainable performance at the same training budget, so mitigation can be interpreted as recovery of damage attributable to the imposed imbalance rather than as an unconstrained gain over a convenient baseline.

The benchmark addresses three research questions:
\begin{description}
  \item[RQ1.] \textbf{Effect of imbalance.} At fixed task, patient split, encoder, training-set size, and evaluation set, how does increasing training imbalance affect tail-class discrimination and calibration in patch classification and WSI-level MIL?
  \item[RQ2.] \textbf{Mitigation effectiveness.} Which mitigation methods recover performance lost specifically to training imbalance, and what are their calibration and computational costs?
  \item[RQ3.] \textbf{Predictors of mitigation headroom.} Do independently measured class separability and effective sample size predict mitigation recovery better than the head-to-tail ratio, and does this differ between patch and WSI supervision?
\end{description}
RQ1 establishes whether a recoverable problem exists. RQ2 conditions claims on that deficit. RQ3 uses measurements made independently of the mitigation models, avoiding the circular use of CE test performance as a separability proxy. The datasets and preprocessing provenance are documented in the companion dataset report, and all detailed method objectives and variants are defined in the companion methods report \citep{queisler2026datasets,queisler2026imbalanceMethods}.

\section{Datasets and Prediction Regimes}
\label{sec:data}
Table~\ref{tab:datasets} lists the diagnostic targets and split units. A patient or case is assigned to exactly one partition before any examples are selected. Where one patient contributes multiple slides, all slides follow the patient assignment. A fixed eligibility pipeline removes unusable samples before splitting and is identical across imbalance conditions.

\begin{table}[ht]
  \centering
  \footnotesize
  \begin{tabularx}{\linewidth}{@{}lXXl@{}}
    \toprule
    Dataset & Patch target & WSI target & Split unit \\
    \midrule
    CAMELYON16 & Mask-derived tumor versus normal & Metastasis versus normal & Slide/patient \\
    TCGA-UT & Cancer type & Cancer type & TCGA participant \\
    BRACS & Seven breast-lesion subtypes from annotated ROIs & Seven breast-lesion subtypes & Patient \\
    PANDA & Mask-derived cancer versus benign & ISUP grade 0--5 & Biopsy \\
    \bottomrule
  \end{tabularx}
  \caption{Prediction targets and leakage-prevention units. PANDA regimes have different targets and are not treated as paired patch-versus-WSI measurements.}
  \label{tab:datasets}
\end{table}

\paragraph{Patch classification.}
Eligible $256\times256$ patches are embedded once with frozen Virchow2. Concatenating the CLS token and mean patch token gives 2560 input features \citep{vorontsov2024virchow,zimmermann2024virchow2}. Patch eligibility, spatial sampling, label derivation, and the deterministic per-slide evidence cap are fixed before condition-specific training subsampling. The common classifier is a two-layer MLP with a 512-unit hidden representation, ReLU, dropout, and a linear class output, except where a method definition requires a different head.

\paragraph{WSI-level MIL.}
A slide is a bag of frozen patch embeddings. A common attention-MIL model maps instances to a 256-dimensional space, assigns normalized attention weights, and classifies the weighted bag representation. The per-slide instance cap and deterministic instance-selection seed are fixed across conditions. Slide labels are the only supervision; patch labels or masks are not used by the primary MIL training runs.

\paragraph{Comparative scope.}
TCGA-UT and BRACS are the primary patch-versus-WSI comparisons because both regimes predict the same labels on related cohorts. PANDA patch and WSI analyses are reported separately because binary cancer detection and ordinal grade-group prediction are different targets. CAMELYON16 is primarily a detection setting and additionally supports a prespecified instance-sparsity analysis using tumor masks.

\subsection{Imbalance and support quantities}
For class supports $N_1,\ldots,N_K$, the headline imbalance ratio is
\begin{equation}
  \rho=\frac{N_{\max}}{N_{\min}}.
  \label{eq:rho}
\end{equation}
The full distribution is also summarized by normalized entropy,
\begin{equation}
  \Hnorm=1-\frac{-\sum_{c=1}^{K}p_c\log p_c}{\log K},
  \qquad p_c=\frac{N_c}{\sum_jN_j}.
  \label{eq:entropy}
\end{equation}
The two quantities are complementary: $\rho$ records the extreme head--tail contrast, whereas $\Hnorm$ uses all classes and is normalized for class count. Both are computed separately at patch and slide level from each realized training manifest. Per-class patient, slide, and patch counts are reported alongside them.

\section{Controlled Experimental Design}
\label{sec:design}
\subsection{Split-first construction and invariant held-out cohorts}
Patient-disjoint training, natural-validation, and natural-test partitions are created first. Imbalance manipulation is confined to training support. Within a split, every condition uses exactly the same validation patients, test patients, and eligible held-out examples. No held-out example is deleted, duplicated, or reweighted according to the training severity. Patch-level evaluation uses every eligible held-out patch under the fixed evidence rule; slide-level evaluation uses every eligible held-out slide. Consequently, recall is estimated from all available held-out cases and prevalence-dependent precision and macro F1 are compared at one fixed test prevalence.

The validation cohort also retains its natural distribution. It is used for hyperparameter and checkpoint selection but never resized to resemble a training condition. Training, validation, and test manifests are versioned with patient identifiers, class counts, construction seed, and content hashes before model fitting.

\subsection{Training conditions}
Each eligible dataset--regime--split has four conditions:
\begin{enumerate}
  \item the full natural training set, used only as a descriptive anchor;
  \item a size-matched balanced training set;
  \item a moderate imbalance condition targeting $\rho=10$; and
  \item a severe imbalance condition targeting $\rho=100$ where feasible.
\end{enumerate}
The balanced, moderate, and severe conditions contain the same total training support $T$: patches for patch classification and slides for MIL. The largest feasible shared $T$ is chosen from training data alone under unique-support and method-feasibility constraints. The natural anchor may contain a different total and is excluded from the imbalance deficit and recovery estimands in Section~\ref{sec:recovery}.

For a $K$-class exponential construction, an assigned class order receives weights
\begin{equation}
  w_i(\rho)=\rho^{-i/(K-1)},\qquad
  n_i(\rho)\approx T\frac{w_i(\rho)}{\sum_{j=0}^{K-1}w_j(\rho)},
  \quad i=0,\ldots,K-1.
  \label{eq:construction}
\end{equation}
Constrained integer allocation preserves $\sum_i n_i=T$, maintains positive support, and never exceeds available unique training examples. Binary targets use $n_{\mathrm{head}}/n_{\mathrm{tail}}=\rho$ with the same total $T$. A prespecified step construction may be used for clinically meaningful BRACS groups; its class grouping and realized supports are reported before analysis. Samples retained at greater severity are nested within the corresponding less-severe training pool where the allocation permits.

If a target ratio is infeasible, the construction uses the largest attainable ratio and reports the requested ratio, achieved ratio, limiting class, and realized counts. It never reduces validation or test support to manufacture feasibility. In particular, TCGA-UT patch classification uses balanced, $\rho=10$, and its maximum feasible or native-severity condition instead of forcing $\rho=100$.

\subsection{Tail assignments and randomization}
Class identity must not be confounded with tail status. Multiclass targets therefore use three locked assignments where feasible: (i) a clinical or native-support assignment, (ii) a reversed assignment for ordinal targets or a deterministic rotation for nominal targets, and (iii) one random class permutation drawn once and recorded before fitting. Binary tasks use both head/tail orientations. If support constraints make an assignment infeasible at the shared $T$, the achieved allocation is reported and the assignment is omitted only under a prespecified feasibility rule.

Three seed families remain independent: patient-split seeds determine cohort membership; assignment seeds determine the locked class order; and initialization seeds determine optimization randomness. Evidence-selection and bootstrap seeds are separately recorded. No seed is chosen based on test performance.

\subsection{Secondary CAMELYON16 MIL experiment}
CAMELYON16 permits separate manipulation of between-slide imbalance and within-slide informative-instance sparsity, two imbalance levels emphasized by SC-MIL \citep{juyal2024scMil}. The number of positive training slides and the fraction of mask-positive instances retained within each tumor bag are crossed factorially. Total training-slide support and bag size remain fixed by replacing removed positive instances with mask-negative instances from the same slide when available. Natural validation and test bags are unchanged. The primary analysis tests main effects and their interaction on WSI balanced accuracy; calibration and attention concentration are secondary outcomes. This analysis is secondary and does not replace the primary controlled conditions.

\section{Mitigation Methods}
\label{sec:methods}
Table~\ref{tab:roster} is intentionally compact. Equations, assumptions, defaults, and fidelity qualifications reside in the methods report \citep{queisler2026imbalanceMethods}. CE and MIL CE are the size-matched no-mitigation references. Balanced Softmax is related prior-correction literature, not an additional tested method \citep{ren2020balancedMeta}.

\begin{table}[ht]
  \centering
  \footnotesize
  \begin{tabularx}{\linewidth}{@{}p{0.25\linewidth}Xl@{}}
    \toprule
    Family & Tested methods & Regime \\
    \midrule
    Reference & CE; MIL CE & Patch; WSI \\
    Sampling & Balanced sampling & Both \\
    Loss and prior correction & Weighted CE; focal loss; train-time and post-hoc logit adjustment; CFAL & Both, except CFAL patch only \\
    Hybrid objective & CE + soft F1; CE + soft MCC; OKO set learning & Patch \\
    Feature mixing & RankMix-inspired feature mixing & WSI \\
    Representation/classifier & cRT; SC-MIL & Both for cRT; WSI for SC-MIL \\
    Dual expert & MDE-inspired dual-expert MIL, including no-consistency ablation & WSI \\
    \bottomrule
  \end{tabularx}
  \caption{Prespecified method roster. ``Both'' means the patch and WSI implementations defined in the companion methods report, not a comparison of their raw scores.}
  \label{tab:roster}
\end{table}

RankMix-inspired denotes the tested adaptation with teacher ranking, representative subsequences, feature interpolation, mixed labels, and a reinitialized soft-label-CE student, but without the source paper's complete combined instance and bag objectives \citep{chen2023rankmix}. MDE-inspired denotes the dual-expert component under frozen Virchow2 bags without multimodal distillation; $\lambda_{\mathrm{con}}=0$ is the dual-expert-without-consistency ablation \citep{ling2025mdeMil}. SC-MIL uses class-aware batches and records valid anchors and positive-pair counts. Logit adjustment uses training-condition priors only \citep{menon2021longTail}; cRT reinitializes and retrains only the final classifier after freezing the learned patch representation or WSI instance encoder and attention aggregator \citep{kang2020decouplingRepresentation}.

\section{Training, Selection, and Replication}
\label{sec:training}
The frozen encoder, classifier family, optimizer family, batch definition, augmentation, numerical precision, and maximum evidence budget are held fixed within a regime. A common maximum epoch or update budget is fixed before the definitive comparison using training and natural-validation learning curves that do not access test labels. Each run stores validation metrics at every checkpoint. The selected checkpoint maximizes natural-validation balanced accuracy, with macro F1 and then lower NLL as deterministic tie-breakers. Test predictions are produced once from that locked checkpoint.

Every tunable mitigation method receives the same number of candidate values for its single imbalance-specific control. The common grids are prespecified in Appendix~\ref{app:controls}. Fixed method parameters follow the definitions in the methods report. cRT records stage-one and stage-two compute separately. Post-hoc logit adjustment has no retraining cost but is selected under the same validation-only rule. Temperature scaling, when reported, fits one positive scalar on natural-validation NLL after model and checkpoint selection and applies it unchanged to test logits \citep{guo2017calibration}.

Controlled contrasts are paired on split, assignment, initialization, and held-out cases. Each locked patient split and tail assignment is repeated with five initialization seeds. Three patient-split seeds are used, subject to dataset feasibility; if a smaller number is required, it is declared before test evaluation. The unit of replication and all missing or failed runs are reported. A failed run is not silently replaced using a test-informed seed.

Computational cost is reported as wall-clock training time, accelerator-hours, peak accelerator memory, total and trainable parameter counts, and effective passes through unique training examples. Costs include all stages required by the method but exclude one-time frozen-feature extraction shared by all methods. Hardware and software environments are recorded with each run.

\section{Outcomes and Statistical Analysis}
\label{sec:analysis}
\subsection{Endpoints and aggregation units}
The primary endpoint is balanced accuracy, equal to macro recall for single-label multiclass prediction. Secondary discrimination outcomes are macro F1; per-class recall and F1; and head, body, and tail recall under the locked assignment. Overall accuracy is descriptive. Patch predictions are summarized in two ways: patch-micro estimates treat eligible patches as predictions while uncertainty remains patient-clustered, and slide/patient-macro estimates first aggregate the metric contribution within each slide or patient so that heavily tiled slides do not dominate. WSI outcomes are slide-level with patient clustering where patients contribute multiple slides.

Probability quality is assessed with NLL, multiclass Brier score, and reliability diagrams. ECE is secondary because it depends on binning and can be unstable with small classwise samples \citep{brier1950verification,naeini2015obtaining}. The number and boundaries of confidence bins are fixed before evaluation, and uncertainty accompanies ECE. Raw and validation-temperature-scaled probabilities are both retained; temperature scaling does not alter discrimination.

\subsection{Imbalance deficit and recovery}
\label{sec:recovery}
For a higher-is-better metric $M$, define the imbalance-induced deficit and recovery fraction at each imbalanced condition as
\begin{equation}
  \begin{aligned}
    D_M &= M(\text{balanced CE})-M(\text{imbalanced CE}),\\
    R_M &= \frac{M(\text{method on imbalanced data})
      -M(\text{imbalanced CE})}{D_M}.
  \end{aligned}
  \label{eq:recovery}
\end{equation}
$R_M=0$ denotes no recovery, $R_M=1$ recovery to the size-matched balanced CE reference, $R_M<0$ further deterioration, and $R_M>1$ performance beyond that reference. The full natural condition is descriptive and never enters $D_M$ or $R_M$.

Mitigation effectiveness is evaluated only for dataset--regime--severity--assignment cells in which the CE balanced-accuracy deficit is at least \num{0.02} and its paired 95\% confidence interval excludes zero. This gate is applied using CE runs before comparing mitigation methods and prevents unstable recovery ratios when there is no evidenced damage to recover. All CE deficits are still reported, including gated-out cells. Recovery for secondary metrics is descriptive and is not used to reopen a cell that fails the primary gate. For lower-is-better calibration outcomes, paired changes are reported in their natural orientation rather than forced into Equation~\eqref{eq:recovery}.

\subsection{Uncertainty and multiplicity}
Confidence intervals use paired hierarchical bootstrap resampling at the highest independent sampling unit: patients, or slides when slide and patient coincide. All patches or slides belonging to a resampled patient travel together. Patches are never bootstrapped as independent observations. Within each bootstrap replicate, paired methods are evaluated on the same resampled held-out cases and matched run seeds. Percentile 95\% intervals are reported with the number of bootstrap replicates and seed fixed in the analysis manifest.

The primary family consists of balanced-accuracy recovery comparisons against imbalanced CE within gated cells. Two-sided paired tests are adjusted by Holm's procedure within each dataset and regime. Effect estimates and intervals remain primary; adjusted $p$-values are supporting evidence. Secondary classwise, calibration, and cost analyses are labelled exploratory unless explicitly prespecified above.

\subsection{Predictors of mitigation headroom}
Separability is measured before mitigation fitting from frozen embeddings using (i) balanced $k$-nearest-neighbour macro recall, (ii) a class-balanced linear probe, and (iii) per-class nearest-neighbour error. Probes use only the controlled training set for fitting or neighbour reference and the unchanged natural-validation set for measurement. They never use test labels or CE test performance. Patch probes operate on patch embeddings. WSI probes use a fixed non-learned summary of the capped instances, such as their mean, so that separability does not inherit the attention-MIL model being explained.

Effective sample size accounts for clustered patch evidence:
\begin{equation}
  N_{\mathrm{eff},c}=\frac{N_c}{1+(\overline m_c-1)\operatorname{ICC}_c},
  \label{eq:neff}
\end{equation}
where $\overline m_c$ is mean cluster size and $\operatorname{ICC}_c$ is the within-patient or within-slide intraclass correlation estimate for class $c$. Slide count itself is used for WSI targets, with patient clustering when necessary.

Recovery is modeled against separability, log imbalance ratio, and log effective sample size. The model includes dataset/task effects, method effects, and patch-versus-WSI interactions; predictors are standardized within regime. The primary comparison asks whether out-of-sample fit improves when independent separability and effective sample size replace or augment $\log\rho$. TCGA-UT and BRACS identify the regime interaction; PANDA is not used as a paired-regime observation because its targets differ. Model form, interaction hierarchy, and cross-validation grouping are fixed before recovery outcomes are inspected.

\section{Design Limitations and Scope}
The protocol isolates training-distribution effects under frozen Virchow2 representations; it does not establish how end-to-end encoder fine-tuning would respond to imbalance. Constructed tails improve causal interpretability but need not reproduce referral pathways or biological prevalence. A shared $T$ controls data volume at the cost of discarding some available training examples, which is why the full natural condition is retained as a descriptive anchor. Some methods change more than one mechanism, so the benchmark supports method-level comparisons rather than universal causal claims about sampling or loss design. The number of independent patients, not the number of patches, limits inference. Finally, calibration on a natural held-out cohort is conditional on that cohort's prevalence and should not be transported to a new deployment population without explicit prior-shift analysis.

\appendix
\section{Protocol Controls}
\label{app:controls}
Table~\ref{tab:controls} records equal-cardinality candidate grids. These values may be narrowed only before definitive test evaluation and with a documented validation-only rationale; they are never changed using test results.

\begin{table}[ht]
  \centering
  \footnotesize
  \begin{tabularx}{\linewidth}{@{}lX@{}}
    \toprule
    Method or family & Candidate imbalance-specific control \\
    \midrule
    Balanced sampling / weighted CE & strength $\in\{0.25,0.5,0.75,1\}$ \\
    Focal loss & $\gamma\in\{0.5,1,1.5,2\}$ \\
    Logit adjustment & $\tau\in\{0.25,0.5,1,2\}$ \\
    CE + soft F1 / soft MCC & auxiliary weight $\in\{0.25,1,4,16\}$ \\
    CFAL & affinity bandwidth $\sigma\in\{0.25,1,2,4\}$ \\
    OKO & odd-class count $k\in\{1,2,4,8\}$, capped by $K-1$ \\
    RankMix-inspired & beta shape $\alpha\in\{0.5,1,2,4\}$ \\
    SC-MIL & temperature $\tau\in\{0.05,0.1,0.2,0.5\}$ \\
    MDE-inspired & consistency weight $\lambda_{\mathrm{con}}\in\{0,0.1,0.25,0.5\}$ \\
    \bottomrule
  \end{tabularx}
  \caption{Prespecified equal-cardinality control grids. cRT has no imbalance-strength grid; its stage-two checkpoint is selected on natural validation.}
  \label{tab:controls}
\end{table}

For every realized condition, the protocol record contains dataset version; eligibility rules; patient, slide, and patch identifiers by partition; requested and achieved $T$ and $\rho$; class assignment; all seed families; evidence caps; model and optimizer configuration; candidate grid; selected checkpoint; environment identifier; and test-prediction hash. Deviations are dated and justified before the affected test labels are accessed.

\section{Terminology}
\begin{description}
  \item[Natural anchor.] The full eligible natural training partition. It is descriptive and may have a different size from controlled conditions.
  \item[Size-matched balanced reference.] A training subset with total support $T$ and approximately equal class counts; it defines the counterfactual reference in $D_M$.
  \item[Imbalance severity.] The requested training allocation summarized by $\rho$; achieved severity is computed from realized integer counts.
  \item[Tail assignment.] The locked mapping from semantic classes to positions in the constructed support profile.
  \item[Imbalance deficit.] The paired balanced-CE minus imbalanced-CE difference at common $T$ and on the same held-out cases.
  \item[Recovery.] A method's improvement over imbalanced CE divided by the evidenced imbalance deficit.
  \item[Patch-micro estimate.] A metric computed over eligible patch predictions, with patient-clustered uncertainty.
  \item[Patient-macro estimate.] An equal-weight aggregation over slides or patients, limiting domination by heavily tiled cases.
  \item[Informative-instance sparsity.] The fraction of instances in a positive bag that contain mask-identified positive tissue.
\end{description}

\clearpage
\bibliographystyle{plainnat}
\bibliography{references}

\end{document}
